% =====================================================================
%  refs.tex —— 交叉引用宏
%
%  为什么不用 cleveref：
%    cleveref 与 ctex 不兼容。最小复现是
%        \documentclass{ctexbook}\usepackage{cleveref}
%        \begin{document}\cref{s1}\par\section{X}\label{s1}\end{document}
%    第二遍编译必报 "Extra \endcsname"。已确认与 hyperref 无关
%    （ctexbook+cleveref 已复现；book+hyperref+cleveref 正常）。
%    zref-clever 与 ctex 兼容，但不提供中文引用名，而本文档要求
%    "第 10 章""式 (3)""表 2"这样的中文形态。
%
%  因此本文档自带一组\emph{按类型区分}的引用宏，建立在 \hyperref + \ref*
%  之上：零宏包风险，输出形态完全可控，且链接照常可用。
%
%  纪律：正文中禁止写死"见 §4.2"这类字面编号，一律使用下列宏。
%        新增引用类型时在此处扩展，不要在正文里就地拼字符串。
% =====================================================================

%  \nref{label}{text} —— 带超链接的文本（\ref* 不产生嵌套链接）
\newcommand{\nref}[2]{\hyperref[#1]{#2}}

%  章：第 10 章
\newcommand{\crefchap}[1]{\nref{#1}{第\,\ref*{#1}\,章}}
%  节：第 10.2 节
\newcommand{\crefsec}[1]{\nref{#1}{第\,\ref*{#1}\,节}}
%  式：式 (3)
\newcommand{\crefeq}[1]{式~\eqref{#1}}
%  表：表 2
\newcommand{\creftab}[1]{\nref{#1}{表~\ref*{#1}}}
%  图：图 1
\newcommand{\creffig}[1]{\nref{#1}{图~\ref*{#1}}}
%  附录：附录 A
\newcommand{\crefapp}[1]{\nref{#1}{附录~\ref*{#1}}}
%  算法/代码块
\newcommand{\creflst}[1]{\nref{#1}{代码~\ref*{#1}}}
